% GENERATED FILE. Edit canonical lessons, then run npm run generate:books.
% canonical-source-hash: fnv1a64:505543b9ed8c9b22
% canonical-lessons: TE-C402-sasanasabha, TE-C402-mantri, TE-C402-adhikaram, TE-C402-palincu, TE-C402-amodincu

\chapter{Governing}
\label{ch:te-governing}
\hlchaptermodality{402}

\begin{chapteropening}
\textbf{By the end of this chapter:} \emph{I can talk about ministers, power and passing laws.}

Complete the last lesson of chapter 402: I can talk about ministers, power and passing laws.
\end{chapteropening}

\section[\texorpdfstring{\te{శాసనసభ} (śāsanasabha)}{శాసనసభ (śāsanasabha)}]{\te{శాసనసభ} (śāsanasabha) — a legislative assembly}

\label{lesson:TE-C402-sasanasabha}



\begin{quote}
Before the new one: say the Telugu for national, then the Telugu for politics.
\end{quote}

\subsection*{You'll want to know: \te{శాసనసభ}}
\textbf{\te{శాసనసభ}} — \emph{śāsanasabha} — \textquotedblleft{}a legislative assembly\textquotedblright{}.

\begin{quote}
\te{కొత్త} \te{చట్టం} \te{గురించి} \te{శాసనసభలో} \te{చర్చ} \te{జరిగింది}.
\end{quote}

\begin{quote}
\emph{kotta caṭṭaṁ guriñci śāsanasabhalō carca jarigindi.}
\end{quote}

\begin{quote}
There was a debate in the assembly about the new law.
\end{quote}

\te{శాసన}, law-making, + \te{సభ}, an assembly. Its members are \te{శాసనసభ్యులు}.

\begin{cousinweb}
\textbf{\te{శాసనసభ}} (\emph{śāsanasabha}) is built from Sanskrit \emph{śāsana}, \textquotedblleft{}instructing; a command, a royal edict\textquotedblright{} (CDIAL 12420), and \emph{sabhā}, \textquotedblleft{}an assembly, a hall\textquotedblright{} (CDIAL 13167).
\end{cousinweb}

\begin{grammarlens}[title={What you've built}]
One more word for giving reasons and opinions.
\end{grammarlens}

\subsection*{Your turn}
\begin{itemize}
\raggedright
  \item \emph{Say it:} \emph{śāsanasabha}
  \item \emph{Say it:} \emph{śāsanasabha}, once more
  \item \emph{Say it:} \emph{kotta caṭṭaṁ guriñci śāsanasabhalō carca jarigindi.} — There was a debate in the assembly about the new law.
  \item \emph{Recall:} say the Telugu for national, then the Telugu for politics, then say \emph{śāsanasabha} again
\end{itemize}

\begin{tcolorbox}[breakable,colback=teal!4,colframe=teal!35!black,title={Before you move on}]
What does \te{శాసనసభ} mean? (\textquotedblleft{}A legislative assembly\textquotedblright{}.) Say it once more.
\end{tcolorbox}
\section[\texorpdfstring{\te{మంత్రి} (mantri)}{మంత్రి (mantri)}]{\te{మంత్రి} (mantri) — a minister}

\label{lesson:TE-C402-mantri}



\begin{quote}
Before the new one: say the Telugu for politics, then the Telugu for a legislative assembly.
\end{quote}

\subsection*{You'll want to know: \te{మంత్రి}}
\textbf{\te{మంత్రి}} — \emph{mantri} — \textquotedblleft{}a minister\textquotedblright{}.

\begin{quote}
\te{విద్యా} \te{మంత్రి} \te{పాఠశాలను} \te{చూడటానికి} \te{వచ్చారు}.
\end{quote}

\begin{quote}
\emph{vidyā mantri pāṭhaśālanu cūḍaṭāniki vaccāru.}
\end{quote}

\begin{quote}
The education minister came to visit the school.
\end{quote}

Plural: \te{మంత్రులు}, \emph{mantrulu}.

\te{ముఖ్యమంత్రి}, the chief minister of a state.

\begin{cousinweb}
\textbf{\te{మంత్రి}} (\emph{mantri}) is borrowed from Sanskrit \emph{mantrin}, \textquotedblleft{}wise, eloquent\textquotedblright{}, and as a noun \textquotedblleft{}a counsellor\textquotedblright{} (CDIAL 9838), from \emph{mantra}, \textquotedblleft{}counsel\textquotedblright{}.
\end{cousinweb}

\begin{grammarlens}[title={What you've built}]
One more word for giving reasons and opinions.
\end{grammarlens}

\subsection*{Your turn}
\begin{itemize}
\raggedright
  \item \emph{Say it:} \emph{mantri}
  \item \emph{Say it:} \emph{mantri}, once more
  \item \emph{Say it:} \emph{vidyā mantri pāṭhaśālanu cūḍaṭāniki vaccāru.} — The education minister came to visit the school.
  \item \emph{Recall:} say the Telugu for politics, then the Telugu for a legislative assembly, then say \emph{mantri} again
\end{itemize}

\begin{tcolorbox}[breakable,colback=teal!4,colframe=teal!35!black,title={Before you move on}]
What does \te{మంత్రి} mean? (\textquotedblleft{}A minister\textquotedblright{}.) Say it once more.
\end{tcolorbox}
\section[\texorpdfstring{\te{అధికారం} (adhikāraṁ)}{అధికారం (adhikāraṁ)}]{\te{అధికారం} (adhikāraṁ) — power; authority}

\label{lesson:TE-C402-adhikaram}



\begin{quote}
Before the new one: say the Telugu for a legislative assembly, then the Telugu for a minister.
\end{quote}

\subsection*{You'll want to know: \te{అధికారం}}
\textbf{\te{అధికారం}} — \emph{adhikāraṁ} — \textquotedblleft{}power; authority\textquotedblright{}.

\begin{quote}
\te{ఆ} \te{పని} \te{చేసే} \te{అధికారం} \te{అతనికి} \te{లేదు}.
\end{quote}

\begin{quote}
\emph{ā pani cēsē adhikāraṁ ataniki lēdu.}
\end{quote}

\begin{quote}
He has no authority to do that.
\end{quote}

\te{అధికారి}, an officer, which you know, is the one who holds it.

\begin{cousinweb}
\textbf{\te{అధికారం}} (\emph{adhikāraṁ}) is borrowed from Sanskrit \emph{adhikāra}, \textquotedblleft{}authority; government, rule, administration\textquotedblright{} (Monier-Williams); \te{అధికారి}, an officer, is the same word.
\end{cousinweb}

\begin{grammarlens}[title={What you've built}]
One more word for giving reasons and opinions.
\end{grammarlens}

\subsection*{Your turn}
\begin{itemize}
\raggedright
  \item \emph{Say it:} \emph{adhikāraṁ}
  \item \emph{Say it:} \emph{adhikāraṁ}, once more
  \item \emph{Say it:} \emph{ā pani cēsē adhikāraṁ ataniki lēdu.} — He has no authority to do that.
  \item \emph{Recall:} say the Telugu for a legislative assembly, then the Telugu for a minister, then say \emph{adhikāraṁ} again
\end{itemize}

\begin{tcolorbox}[breakable,colback=teal!4,colframe=teal!35!black,title={Before you move on}]
What does \te{అధికారం} mean? (\textquotedblleft{}Power; authority\textquotedblright{}.) Say it once more.
\end{tcolorbox}
\section[\texorpdfstring{\te{పాలించు} (pāliñcu)}{పాలించు (pāliñcu)}]{\te{పాలించు} (pāliñcu) — to govern; to rule}

\label{lesson:TE-C402-palincu}



\begin{quote}
Before the new one: say the Telugu for a minister, then the Telugu for power.
\end{quote}

\subsection*{You'll want to know: \te{పాలించు}}
\textbf{\te{పాలించు}} — \emph{pāliñcu} — \textquotedblleft{}to govern; to rule\textquotedblright{}.

\begin{quote}
\te{ఆ} \te{కుటుంబం} \te{ఆ} \te{ప్రాంతాన్ని} \te{చాలా} \te{కాలం} \te{పాలించింది}.
\end{quote}

\begin{quote}
\emph{ā kuṭumbaṁ ā prāntānni cālā kālaṁ pāliñcindi.}
\end{quote}

\begin{quote}
That family ruled the region for a long time.
\end{quote}

Past (I …): \te{పాలించాను}, \emph{pāliñcānu}.

The noun is \te{పాలన}, rule.

\begin{cousinweb}
\textbf{\te{పాలించు}} (\emph{pāliñcu}) is built on Sanskrit \emph{pālana}, \textquotedblleft{}protecting, nourishing\textquotedblright{} (CDIAL 8128), with the Telugu verb ending -\te{ఇంచు}; Telugu uses it for ruling and governing.
\end{cousinweb}

\begin{grammarlens}[title={What you've built}]
One more word for giving reasons and opinions.
\end{grammarlens}

\subsection*{Your turn}
\begin{itemize}
\raggedright
  \item \emph{Say it:} \emph{pāliñcu}
  \item \emph{Say it:} \emph{pāliñcu}, once more
  \item \emph{Say it:} \emph{ā kuṭumbaṁ ā prāntānni cālā kālaṁ pāliñcindi.} — That family ruled the region for a long time.
  \item \emph{Recall:} say the Telugu for a minister, then the Telugu for power, then say \emph{pāliñcu} again
\end{itemize}

\begin{tcolorbox}[breakable,colback=teal!4,colframe=teal!35!black,title={Before you move on}]
What does \te{పాలించు} mean? (\textquotedblleft{}To govern; to rule\textquotedblright{}.) Say it once more.
\end{tcolorbox}
\section[\texorpdfstring{\te{ఆమోదించు} (āmōdiñcu)}{ఆమోదించు (āmōdiñcu)}]{\te{ఆమోదించు} (āmōdiñcu) — to approve; to pass (a bill)}

\label{lesson:TE-C402-amodincu}



\begin{quote}
Before the new one: say the Telugu for power, then the Telugu for to govern.
\end{quote}

\subsection*{You'll want to know: \te{ఆమోదించు}}
\textbf{\te{ఆమోదించు}} — \emph{āmōdiñcu} — \textquotedblleft{}to approve; to pass (a bill)\textquotedblright{}.

\begin{quote}
\te{శాసనసభ} \te{కొత్త} \te{చట్టాన్ని} \te{ఆమోదించింది}.
\end{quote}

\begin{quote}
\emph{śāsanasabha kotta caṭṭānni āmōdiñcindi.}
\end{quote}

\begin{quote}
The assembly passed the new law.
\end{quote}

Past (I …): \te{ఆమోదించాను}, \emph{āmōdiñcānu}.

The noun is \te{ఆమోదం}, approval.

\begin{cousinweb}
\textbf{\te{ఆమోదించు}} (\emph{āmōdiñcu}) is built on Sanskrit \emph{āmoda}, \textquotedblleft{}joy, serenity, pleasure\textquotedblright{} (Monier-Williams), with the Telugu verb ending -\te{ఇంచు}; Telugu uses it for giving approval.
\end{cousinweb}

\begin{grammarlens}[title={What you've built}]
One more word for giving reasons and opinions.
\end{grammarlens}

\subsection*{Your turn}
\begin{itemize}
\raggedright
  \item \emph{Say it:} \emph{āmōdiñcu}
  \item \emph{Say it:} \emph{āmōdiñcu}, once more
  \item \emph{Say it:} \emph{śāsanasabha kotta caṭṭānni āmōdiñcindi.} — The assembly passed the new law.
  \item \emph{Recall:} say the Telugu for power, then the Telugu for to govern, then say \emph{āmōdiñcu} again
\end{itemize}

\begin{tcolorbox}[breakable,colback=teal!4,colframe=teal!35!black,title={Before you move on}]
What does \te{ఆమోదించు} mean? (\textquotedblleft{}To approve; to pass (a bill)\textquotedblright{}.) Say it once more.
\end{tcolorbox}
